\chapter*{Заключение}
\addcontentsline{toc}{chapter}{Заключение}
\label{ch:conclusion}

В ходе выполнения лабораторной работы~\cfgLabNumber{} был развёрнут
и настроен почтовый сервер \textbf{iRedMail~1.6.2} на виртуальной
машине \texttt{mail} (Ubuntu~22.04~Server, IP~\texttt{192.168.\cfgLabStudentN.5})
в рамках уже существующей инфраструктуры учебного стенда.

В процессе работы были выполнены следующие задачи:
\begin{enumerate}
  \item Создана и настроена виртуальная машина \texttt{mail} со статическим
        IP-адресом и корректным FQDN-именем\\
        \texttt{mail.yazikov.iks531.local}.
  \item В DNS-зону шлюза \texttt{gateway} добавлена A-запись для
        нового почтового сервера.
  \item Загружен и запущен интерактивный установщик iRedMail;
        выбраны Nginx в качестве веб-сервера и OpenLDAP
        для хранения учётных записей.
  \item После завершения установки проверена работа сервисов
        \texttt{postfix}, \texttt{dovecot} и \texttt{nginx}.
  \item Через браузер Desktop подтверждена доступность веб-интерфейсов
        \texttt{iRedAdmin} и \texttt{Roundcube}.
\end{enumerate}

Таким образом, в учебной сети организован полноценный почтовый сервер
с поддержкой SMTP, IMAP, антиспам-фильтрацией и удобным
веб-интерфейсом для пользователей и администратора.

\endinput
